% =====================================================================
%  sample.tex : レイアウト見本（すべての環境を一度ずつ使う）
%  スタイルを変えたら、これをビルドして見た目を確かめる：
%    cd tex && latexmk sample.tex
% =====================================================================
% =====================================================================
%  preamble.tex : main.tex と sample.tex で共通のプリアンブル
% =====================================================================
\documentclass[
  book,                 % 片面ではなく見開き（章は奇数ページから）
  paper=b5j,            % JIS B5（182mm × 257mm）
  fontsize=10pt,
  line_length=40zw,     % 1行 40 字
  number_of_lines=34,   % 1ページ 34 行
  gutter=22mm,          % のど（とじ側）の余白
  head_space=24mm,      % 天（上）の余白
  baselineskip=1.7zw,
]{jlreq}

\usepackage{style/textbook}


\begin{document}

\part{統計基礎}

\chapter{確率と確率分布}
\label{chap:sample}

\begin{chaptergoal}
  \begin{itemize}
    \item 積率母関数の定義と使い方を説明できる。
    \item チェビシェフの不等式を導き、その意味を言葉で説明できる。
  \end{itemize}
\end{chaptergoal}

\section{確率分布、確率変数}
\label{sec:sample-moments}

\keywords{チェビシェフの不等式、積率、尖度、歪度、積率母関数}

\subsection{積率母関数}

確率変数の性質を調べるとき、期待値 $\E[X]$ や分散 $\V[X]$ を1つずつ計算するのは手間がかかります。
そこで、すべての積率をまとめて持っている関数を考えます。

\begin{definition}[積率母関数][def:mgf]
  確率変数 $X$ に対して、
  \begin{equation}
    M_X(t) = \E\left[e^{tX}\right]
    \label{eq:mgf}
  \end{equation}
  を $X$ の\term{積率母関数}{せきりつぼかんすう}といいます。
\end{definition}

\begin{intuition}
  $e^{tX}$ を展開すると $1 + tX + \frac{t^2}{2!}X^2 + \cdots$ となります。
  つまり $M_X(t)$ は、$\E[X], \E[X^2], \dots$ を「係数として1列に並べて保存した入れ物」です。
\end{intuition}

\begin{theorem}[積率母関数から積率を取り出す][thm:mgf-moment]
  $M_X(t)$ を $k$ 回微分して $t=0$ を代入すると、$k$ 次の積率が得られる：
  \[
    M_X^{(k)}(0) = \E[X^k].
  \]
\end{theorem}

\begin{proof}
  \cref{eq:mgf} の両辺を $t$ で $k$ 回微分すると $M_X^{(k)}(t) = \E[X^k e^{tX}]$ です。
  ここで $t=0$ とおけば $e^{0}=1$ なので、$M_X^{(k)}(0)=\E[X^k]$ となります。
\end{proof}

\begin{merit}
  積分（または和）の計算を1回すませれば、あとは\textbf{微分するだけで}平均も分散も求まります。
  微分は積分よりずっと楽な計算です。
\end{merit}

\begin{mathnote}[$e^x$ のテイラー展開]
  $e^x = 1 + x + \dfrac{x^2}{2!} + \dfrac{x^3}{3!} + \cdots$ が成り立ちます。
  くわしくは数学基礎「1変数関数の微分法」を参照してください。
\end{mathnote}

\begin{example}[ポアソン分布の積率母関数][ex:poisson-mgf]
  $X \sim \Po(\lambda)$ のとき、
  \[
    M_X(t) = \sum_{x=0}^{\infty} e^{tx}\,\frac{\lambda^x e^{-\lambda}}{x!}
           = e^{-\lambda}\sum_{x=0}^{\infty}\frac{(\lambda e^t)^x}{x!}
           = \exp\bigl(\lambda(e^t-1)\bigr).
  \]
\end{example}

\begin{problem}[ポアソン分布の平均]
  \cref{ex:poisson-mgf} の結果を使って、$\E[X]$ を求めなさい。
  \tcblower
  $M_X'(t) = \lambda e^t \exp\bigl(\lambda(e^t-1)\bigr)$ なので、
  \cref{thm:mgf-moment} より $\E[X] = M_X'(0) = \lambda$ です。
\end{problem}

\begin{supplement}[確率変数の変換]
  $Y=g(X)$ の確率密度関数は、$g$ が単調なら $f_Y(y) = f_X\bigl(g^{-1}(y)\bigr)\abs*{\diff{}{y}g^{-1}(y)}$ で求まります。
\end{supplement}

\begin{reference}[ジャックナイフ法]
  データを1つずつ除いて推定をくり返す方法です。名前だけ知っておけば十分です。
\end{reference}

\begin{exampoint}
  「歪度は3次、尖度は4次の標準化積率」という対応は、そのまま問われます。
\end{exampoint}

\subsection{図とコード}

\Cref{fig:sample-pgf} は TikZ（pgfplots）で、\cref{fig:sample-py} は Python で描いた図です。

\begin{figure}[tb]
  \centering
  \begin{tikzpicture}
    \begin{axis}[width=0.7\linewidth, height=4.5cm, axis lines=left,
                 xlabel={$x$}, ylabel={$f(x)$}, ymin=0, samples=100, domain=-4:4]
      \addplot[thick, tbMain] {exp(-x^2/2)/sqrt(2*pi)};
      \addplot[thick, tbTheorem, dashed] {exp(-abs(x))/2};
    \end{axis}
  \end{tikzpicture}
  \caption{標準正規分布（実線）とラプラス分布（破線）}
  \label{fig:sample-pgf}
\end{figure}

\begin{figure}[tb]
  \centering
  \includegraphics{sample_normal.pdf}
  \caption{Python（matplotlib）で描いた図の例}
  \label{fig:sample-py}
\end{figure}

\begin{pycode}[正規乱数の平均]
import numpy as np

rng = np.random.default_rng(seed=0)  # シードを固定して再現できるようにする
x = rng.normal(loc=0.0, scale=1.0, size=1000)
print(x.mean())
\end{pycode}

\begin{pyoutput}
0.04440216000254035
\end{pyoutput}

\begin{table}[tb]
  \centering
  \caption{表の例}
  \label{tab:sample}
  \begin{tabular}{lcc}
    \toprule
    分布 & 期待値 & 分散 \\
    \midrule
    $\Bin(n,p)$ & $np$ & $np(1-p)$ \\
    $\Po(\lambda)$ & $\lambda$ & $\lambda$ \\
    $\Normal(\mu,\sigma^2)$ & $\mu$ & $\sigma^2$ \\
    \bottomrule
  \end{tabular}
\end{table}

\Cref{tab:sample} のように、表は \texttt{booktabs} の罫線で書きます。
参照の見本：\cref{chap:sample}、\cref{sec:sample-moments}、\cref{def:mgf}、\cref{thm:mgf-moment}。

\begin{chaptersummary}
  \begin{itemize}
    \item 積率母関数 $M_X(t)=\E[e^{tX}]$ を微分すると積率が得られる。
  \end{itemize}
\end{chaptersummary}

\printindex

\end{document}
